% Tableau comparatif de la Phase 0 (Baseline) pour Berlin52
\begin{table}[h]
\centering
\begin{tabular}{|l|c|c|c|}
\hline
\textbf{Algorithme} & \textbf{Distance} & \textbf{Ecart (Gap \%)} & \textbf{Temps (ms)} \\ \hline
Nearest Neighbor (Départ 0) & 8981 & +19,08\% & 5,04 ms \\ \hline
Best Nearest Neighbor (Tous départs) & 8182 & +8,49\% & 181,49 ms \\ \hline
2-opt Séquentiel (Baseline) & 7670 & +1,69\% & 96,71 ms \\ \hline
\textbf{2-opt Lot + K-NN (Phase 1\&2)} & \textbf{7841} & \textbf{+3,97\%} & \textbf{42,81 ms} \\ \hline
\end{tabular}
\caption{Performance des algorithmes de base sur l'instance Berlin52 (Optimum : 7542)}
\label{tab:baseline_berlin52}
\end{table}

% --- JUSTIFICATIONS ET ANALYSE POUR LE RAPPORT ---
% 1. EFFICACITÉ DE LA RECHERCHE LOCALE : Le passage du Best NN au 2nd-opt montre un gain significatif 
%    de qualité (8182 -> 7670), prouvant que la recherche locale est indispensable.
% 2. COMPROMIS VITESSE/PRÉCISION (K-NN) : La version Lot + K-NN est plus rapide (42ms), mais un peu 
%    moins précise (+3,97%) que le 2-opt complet. C'est le prix de la réduction spatiale.
% 3. SPEEDUP ALGORITHMIQUE : Sur ch150, l'optimisation Batch + K-NN offre un Speedup de x12,1 
%    par rapport au séquentiel (0,29s contre 3,51s).
% 4. TRANSITION : Cette accélération massive justifie l'intérêt du parallélisme et du K-NN 
%    qui seront exploités dans les phases de production suivantes.

justifications : 
 passer de 232ms (Best NN) à 94ms (2-opt) montre que la recherche locale est efficace, mais que sur une 
 instance de 1000 villes, ces millisecondes vont devenir des secondes entières, justifiant ainsi tes phases 
 de parallélisation (Phase 1, 2, 3). 